\section{Normalized approximation with independent scales}

\begin{theorem}[An improving family at independent scales]
    \label{thm:independent_long_middle_actual_dg_family}
    \lean{TNLean.MPS.AreaLaw.exists_independent_long_middle_actual_dg_family}
    \leanok
    Let $m\ge q\ge2$, $\ell\ge1$, $t\ge21$, $0<g\le1$ and $0<\varepsilon\le1$.
    Suppose
    \begin{align}
        288e^{11/8}e^{-(t-g)/8} & \le\varepsilon g/10^6,
        \label{eq:independent_family_tail}
    \end{align}
    and suppose that a real number $\Delta$ bounds
    $(2e^{-2\ell^2\sqrt{(4g/5)/R}})^2$ for every $0<R\le3t+4\ell$.
    Put $D=(m\ell^2)^{22\ell}$.
    Let $J=[p,r)$ have length greater than $4\ell$, and let $H$ be an
    original positive unit nearest-neighbour Hamiltonian with full ground
    space $G$ at energy $\mu$ and spectrum in $\{\mu\}\cup[\mu+g,\infty)$.
    There are a target $T$ and $D^2$ operators $F_i$ on $\mathcal H_J$,
    chosen before every viable subspace below, such that $\dim T=\dim G$
    and $T,G$ are mutually $\varepsilon/200$-close.
    If $S\subseteq\mathcal H_J$ is $\eta$-viable for $T$ and $\eta<1$, then
    \begin{align}
        V=\operatorname{span}_i F_iS,
        \qquad \dim V\le D^2\dim S,
        \qquad V\text{ is }\frac{\Delta}{(1-\eta)^2}\text{-viable for }T.
        \label{eq:independent_family_improvement}
    \end{align}
\end{theorem}
\begin{proof}
    \leanok
    \uses{thm:inward_effective_cuts,thm:five_region_scaled_low_cluster,
        thm:actual_clipped_polynomial_family}
    Move the cuts inward to a retained interval whose omitted collar has
    at most $2\ell$ sites. The physical clipping theorem gives the actual
    low cluster and norm budget $3t+4\ell$.
    Construct the common polynomial family at degree $\ell^2$, transport
    it to $J$, and absorb the collar count into $D^2$.
    The actual-cluster Chebyshev estimate gives
    (\ref{eq:independent_family_improvement}); the displayed numerical
    hypothesis bounds its squared error. Extend by zero to exactly $D^2$
    operators without changing their image span.
\end{proof}

\begin{theorem}[The squared filter error]
    \label{thm:squared_chebyshev_error_eq}
    \lean{Real.squared_chebyshev_error_eq}
    \leanok
    For all real $\ell,\eta,R$,
    $(2e^{-2\ell^2\sqrt{\eta/R}})^2=4e^{-4\ell^2\sqrt{\eta/R}}$.
\end{theorem}
\begin{proof}
    \leanok
    Apply $(e^x)^2=e^{2x}$ and expand the scalar square.
\end{proof}

\begin{theorem}[An error proportional to the requested accuracy]
    \label{thm:squared_chebyshev_error_le_accuracy}
    \lean{Real.squared_chebyshev_error_le_accuracy}
    \leanok
    Let $\delta>0$, $a\ge16$, $\log(1/\delta)\le a$ and
    $a\le\ell^2\sqrt{\eta/R}$. Then
    $(2e^{-2\ell^2\sqrt{\eta/R}})^2\le\delta/8$.
\end{theorem}
\begin{proof}
    \leanok
    \uses{thm:squared_chebyshev_error_eq}
    One has $e^{-a}\le\delta$ and $e^{3a}\ge1+3a\ge49$.
    Hence $4e^{-4a}=(4e^{-3a})e^{-a}\le\delta/8$.
\end{proof}

\begin{theorem}[One scalar choice for physical approximation]
    \label{thm:sharp_approximation_scalar_parameters}
    \lean{Real.exists_sharp_approximation_scalar_parameters}
    \leanok
    There is a universal constant $C\ge21$ such that, for every
    $0<g\le1$ and $0<\delta\le1$, putting
    $a=16+\log(1/g)+\log(1/\delta)$ and $t=Ca$, there is an integer
    $\ell\ge1$ with
    \begin{align}
        \ell & \le\frac{8}{4g/5}+(7C)^{1/4}(4g/5)^{-1/4}a^{3/4},
        \label{eq:sharp_approximation_window}
    \end{align}
    such that $t\ge21$, $288e^{11/8}e^{-(t-g)/8}\le\delta g/10^6$,
    and $(2e^{-2\ell^2\sqrt{(4g/5)/R}})^2\le\delta/8$ for every
    $0<R\le3t+4\ell$.
\end{theorem}
\begin{proof}
    \leanok
    \uses{thm:independent_truncation_threshold_constant,
        thm:quarter_power_truncation_filter_scale,
        thm:squared_chebyshev_error_le_accuracy}
    Use the universal threshold constant. Increasing its logarithmic
    variable from $1+\log(1/g)+\log(1/\delta)$ to $a$ decreases the tail.
    Apply the quarter-power window choice with gap $4g/5$ and exponent $a$,
    then the preceding squared-error estimate.
\end{proof}

\begin{theorem}[The logarithm of the independent family count]
    \label{thm:sharp_polynomial_family_log_count}
    \lean{Real.sharp_polynomial_family_log_count}
    \leanok
    Let $q,\ell\ge1$, $0<g\le1$, $a\ge1$, $b\ge0$ and
    $\ell\le10/g+bg^{-1/4}a^{3/4}$.
    Put $u=1+\log q+\log(1/g)$ and $D_1=(q\ell^2)^{22\ell}$.
    Then
    \begin{align}
        \log D_1\le44(1+\log(11+b))
        (10/g+bg^{-1/4}a^{3/4})u(1+\log a).
        \label{eq:sharp_polynomial_log_count}
    \end{align}
\end{theorem}
\begin{proof}
    \leanok
    The window bound implies $\ell\le(11+b)a/g$, so
    $\log\ell\le\log(11+b)+\log(1/g)+\log a$.
    Thus $\log q+2\log\ell\le2(1+\log(11+b))u(1+\log a)$.
    Multiply by $22\ell$ and apply the sharper window bound once more.
\end{proof}

\begin{theorem}[A rank budget with one ground-space factor]
    \label{thm:sharp_approximation_log_rank_budget}
    \lean{Real.sharp_approximation_log_rank_budget}
    \leanok
    Under the hypotheses of Theorem~\ref{thm:sharp_polynomial_family_log_count},
    let $r,D_0\ge1$, $K_0\ge0$ and $\log D_0\le K_0u^3/g$.
    Then
    \begin{align}
        \log(rD_0^{24}D_1^2)
         & \le\log r+
        \bigl(24K_0+88(1+\log(11+b))(11+b)\bigr)
        \bigl(u^3/g+g^{-1/4}a^{3/4}u\bigr)(1+\log a).
        \label{eq:sharp_log_rank_budget}
    \end{align}
\end{theorem}
\begin{proof}
    \leanok
    \uses{thm:sharp_polynomial_family_log_count}
    Expand the left side as $\log r+24\log D_0+2\log D_1$.
    Since $u\ge1$, one has $u/g\le u^3/g$.
    Apply the two logarithmic count bounds and use $1+\log a\ge1$.
\end{proof}

\begin{theorem}[Normalized physical approximation with an integer rank budget]
    \label{thm:physical_dg_normalized_approximation_independent}
    \lean{TNLean.MPS.AreaLaw.exists_physical_dg_normalized_approximation_of_independent_scales}
    \leanok
    There are universal constants $A\ge1$ and $C\ge21$ with the following
    property. Let $H,q,\mu,\nu,\gamma,G,J$ satisfy the physical hypotheses
    of Theorem~\ref{thm:physical_interval_bootstrap}, and put
    \begin{align}
        g      & =\min(\gamma,1),      & u   & =1+\log q+\log(1/g),\notag \\
        \ell_0 & =\lceil Au^2/g\rceil, & D_0 & =(q\ell_0^2)^{22\ell_0}.
    \end{align}
    Then $D_0\ge2$ and
    $\log D_0\le22(A+1)(6+2\log(A+1))u^3/g$.
    For every $0<\delta\le1$, put $a=16+\log(1/g)+\log(1/\delta)$.
    There is an integer $\ell_1\ge1$ satisfying
    (\ref{eq:sharp_approximation_window}), chosen before the vector below.
    Put $D_1=(q\ell_1^2)^{22\ell_1}$.
    Every unit $\psi\in G$ has a unit approximant $\phi$ in the original
    Hilbert space with
    $|\langle\phi,\psi\rangle|\ge1-\delta$ and Schmidt rank at most
    $(\dim G)D_0^{24}D_1^2$ across $J,J^c$.
\end{theorem}
\begin{proof}
    \leanok
    \uses{thm:physical_interval_all_accuracy_family,
        thm:sharp_approximation_scalar_parameters,
        thm:independent_long_middle_actual_dg_family}
    The bootstrap gives a fixed $0.015$-viable $W_0$ with
    $\dim W_0<(\dim G)D_0^{24}$ before the accuracy or vector is chosen.
    If $|J|\le4\ell_1$, use $\phi=\psi$; the full regional dimension
    is at most $D_1^2$.
    Otherwise apply the independent family. Transfer $W_0$ to its target
    with viability error $0.04$, improve to error at most $\delta/5$,
    and transfer back to $G$ with error at most $\delta/2$.
    The enlarged space has dimension at most $D_1^2\dim W_0$.
    Its normalized projection witness, transported back by the actual
    cut isometry, has squared overlap at least $1-\delta/2$ and the stated
    rank bound. Since $0<\delta\le1$, its absolute overlap is at least
    $1-\delta$.
\end{proof}

\begin{theorem}[Normalized approximation with the sharp logarithmic rank bound]
    \label{thm:physical_dg_interval_normalized_approximation}
    \lean{TNLean.MPS.AreaLaw.exists_physical_dg_interval_normalized_approximation}
    \leanok
    There is a universal constant $K>0$ with the following property.
    Let $H,q,\mu,\nu,\gamma,G,J$ satisfy the physical hypotheses of
    Theorem~\ref{thm:physical_interval_bootstrap}.
    Set $r=\dim G$, $g=\min(\gamma,1)$ and $u=1+\log q+\log(1/g)$.
    For every $0<\delta\le1$, put $a=16+\log(1/g)+\log(1/\delta)$.
    Every unit $\psi\in G$ has a unit vector $\phi$ in the original
    Hilbert space such that $|\langle\phi,\psi\rangle|\ge1-\delta$ and
    \begin{align}
        \log\operatorname{SR}_{J:J^c}(\phi)
         & \le\log r+K\bigl(u^3/g+g^{-1/4}a^{3/4}u\bigr)(1+\log a).
        \label{eq:physical_sharp_log_rank}
    \end{align}
    Taking $J$ to be a prefix gives the fixed-cut conclusion of
    \cite[Theorem~4.9]{Arad2016RigorousRG} for $0<\gamma\le1$,
    with explicit logarithmic factors.
    For larger supplied gaps, this assertion uses $g=1$.
    The normalization and the independent parameter choice are recorded in
    \cite{gap:alvv16_approximation_parameters}.
\end{theorem}
\begin{proof}
    \leanok
    \uses{thm:physical_dg_normalized_approximation_independent,
        thm:sharp_approximation_log_rank_budget}
    Use the integer rank budget and set
    $K_0=22(A+1)(6+2\log(A+1))$ and
    $b=(7C)^{1/4}(5/4)^{1/4}$.
    The window bound becomes $\ell_1\le10/g+bg^{-1/4}a^{3/4}$.
    Apply (\ref{eq:sharp_log_rank_budget}) and monotonicity of the logarithm
    to the actual Schmidt rank of the normalized vector.
\end{proof}

\begin{theorem}[One normalized matrix product state at all cuts]
    \label{thm:physical_dg_normalized_mps_approximation}
    \lean{TNLean.MPS.AreaLaw.exists_physical_dg_normalized_isOpenMPS_approximation}
    \leanok
    There is a universal constant $K>0$ with the following property.
    Let $H,q,\mu,\nu,\gamma,G$ satisfy the physical hypotheses of
    Theorem~\ref{thm:physical_interval_bootstrap} on a chain of $N$ sites.
    Put $r=\dim G$, $g=\min(\gamma,1)$ and $u=1+\log q+\log(1/g)$.
    For every $0<\delta\le1$, set
    \begin{align}
        b=16+\log(1/g)+\log\bigl(16(N+1)/\delta\bigr).
        \label{eq:physical_global_accuracy}
    \end{align}
    There is an integer $B\ge1$, chosen before the ground vector, with
    \begin{align}
        \log B\le K\bigl(u^3/g+g^{-1/4}b^{3/4}u\bigr)(1+\log b).
        \label{eq:physical_global_bond}
    \end{align}
    Every unit $\psi\in G$ has a unit open-boundary matrix product state
    $\Phi$ of bond dimension at most $rB$ such that
    $|\langle\Phi,\psi\rangle|\ge1-\delta$ and
    $\|\Phi-\psi\|^2\le\delta/2$.
    This gives the global approximation conclusion of
    \cite[Theorem~4.9]{Arad2016RigorousRG} for $0<\gamma\le1$, with
    explicit logarithmic factors. The capped gap for larger supplied gaps
    is as described in \cite{gap:alvv16_approximation_parameters}.
\end{theorem}
\begin{proof}
    \leanok
    \uses{thm:physical_dg_interval_normalized_approximation,
        thm:chain_sequential_truncation_eigen,
        thm:chain_sequential_truncation_sq}
    Set $\delta_c=\delta/(16(N+1))$.
    The fixed-cut bound gives a common logarithmic rank budget
    $\log r+E$ for every prefix cut, where
    $E=K_0(u^3/g+g^{-1/4}b^{3/4}u)(1+\log b)$ and $K_0$ is universal.
    Choose $B=\lceil e^E\rceil$ and $D=rB$.
    At cut $k$, the normalized approximant of rank at most $D$ has overlap
    at least $1-\delta_c$. The maximal Schmidt-overlap bound therefore gives
    \begin{align}
        \epsilon_k(D)=1-\sum_{i=1}^D\lambda_i(\rho_k)
        \le1-(1-\delta_c)^2\le2\delta_c,
        \label{eq:physical_global_prefix_tail}
    \end{align}
    where $\rho_k$ is the original state's prefix marginal and eigenvalues
    beyond its dimension contribute zero.
    Let $\Pi_k$ project onto its leading $D$ eigenvectors, and define one vector
    \begin{align}
        \Phi_0=(\Pi_1\otimes\Id)\cdots(\Pi_{N-1}\otimes\Id)\psi.
        \label{eq:physical_global_projection_product}
    \end{align}
    The sequential-truncation theorem makes $\Phi_0$ an open-boundary matrix
    product state of bond dimension at most $D$, with
    $\|\Phi_0-\psi\|^2\le\sum_k\epsilon_k(D)\le\delta/8$.
    Thus $\Phi_0\ne0$. Put $\Phi=\Phi_0/\|\Phi_0\|$.
    Normalization at most doubles the norm error, so
    $\|\Phi-\psi\|^2\le\delta/2$ and
    $|\langle\Phi,\psi\rangle|\ge1-\delta$.
    Scalar multiplication preserves the same bond bound.
    Finally $\log B\le E+1$; the extra one is absorbed into the universal
    constant in (\ref{eq:physical_global_bond}).
\end{proof}
